\begin{table}[tp]
\centering
\caption{Gradient sliding for \href{https://github.com/k3nfalt/longtails/blob/main/LongTails/Setups/SlidingGeneralizedSmooth.lean}{SlidingGeneralizedSmooth}.}\label{tab:SlidingGeneralizedSmooth}
\parbox{\linewidth}{\footnotesize Minimize $F=f+h$ over convex $Q$ in a real Hilbert space. The convex differentiable $f$ has continuous gradient and satisfies $\|\nabla f(x)-\nabla f(y)\|\le [L_0+L_1\sup_{z\in[x,y]}\|\nabla f(z)\|]\|x-y\|$ on the whole space. Oracles supply exact gradients, exact projections onto $Q$, and possibly adversarial vectors $s(u)$ with $\|s(u)\|\le M$ and $h(v)\ge h(u)+\langle s(u),v-u\rangle$ for $u,v\in Q$. There is $x^\star\in\arg\min_Q F$ with $x_0\in Q$, $\|x_0-x^\star\|\le R$, $F(x_0)-F(x^\star)\le\Delta$, and $\|\nabla f(x^\star)\|\le G$. Here $L_0,R,\Delta,\varepsilon>0$ and $L_1,M,G\ge0$. Target: a feasible output $x$ with $F(x)-F(x^\star)\le\varepsilon$ almost surely.}\par\smallskip
\begin{threeparttable}
\footnotesize
\setlength{\tabcolsep}{4pt}\renewcommand{\arraystretch}{1.25}
\begin{tabularx}{\linewidth}{@{}>{\raggedright\arraybackslash}p{0.20\linewidth}>{\raggedright\arraybackslash}p{0.32\linewidth}>{\raggedright\arraybackslash}p{0.25\linewidth}>{\raggedright\arraybackslash}X@{}}
\toprule
\textbf{Result} & \textbf{Oracle complexity} & \textbf{Novelty} & \textbf{Limitations} \\
\midrule
\textbf{(1)} \href{https://github.com/k3nfalt/longtails/blob/main/LongTails/Theorems/GeneralizedGradientSliding.lean}{Generalized\allowbreak{}Gradient\allowbreak{}Sliding} \textbf{(Ours)} $(\star\star\star)$ \href{https://github.com/k3nfalt/longtails/blob/main/scripts/reviews/GeneralizedGradientSliding.md}{(review)}
& Total: $\mathcal{O}\!\left(\begin{gathered}\sqrt{AR^2/\varepsilon}+(1+L_1R)\ell\\{}+M^2R^2/\varepsilon^2\end{gathered}\right)$.
\par Gradients: $\mathcal{O}\!\left(\sqrt{\frac{AR^2}{\varepsilon}}+(1+L_1R)\ell\right)$.\par $G$ is used only in the gradient-count cap.
& & \\
\midrule[1.2pt]
\textbf{(2)} Gradient sliding \citep{sgsLan2016}, Cor.~1(a), arXiv v2. Derived bound: finite-dimensional localization to $X=Q\cap\overline B(x_0,2R)$.
& Total: $\mathcal{O}\!\left(1+\sqrt{\frac{Ae^{3L_1R}R^2}{\varepsilon}}+\frac{M^2R^2}{\varepsilon^2}\right)$.
\par Gradients: $\mathcal{O}\!\left(1+\sqrt{\frac{Ae^{3L_1R}R^2}{\varepsilon}}\right)$.
\par Uses $L=Ae^{3L_1R}$ and exact $\Pi_X$; at $L_1=0$, use $X=Q$, $L=L_0$.
& \textbf{(1)}: gradient saving $e^{3L_1R/2}$ when $L_1R\gg1$ and $\sqrt{AR^2/\varepsilon}\gg(1+L_1R)\ell$; total saving additionally requires $M^2R^2/\varepsilon^2\ll\sqrt{AR^2/\varepsilon}$. Only $\Pi_Q$ needed. Matches leading rate for $L_1R\lesssim1$; scope extends global smoothness.
& \textbf{(1)}: extra $\ell$ term at $L_1=0$; no Bregman geometry or additional simple prox term. \\
\midrule
\textbf{(3)} Normalized subgradient \citep{sgsGrimmer2019}, Thm.~2, arXiv v3. Derived bound, $E=\mathbb R^d$; best iterate selected using $F$ values.
& Total and gradients:
$\mathcal{O}\!\left(\begin{gathered}1+(G+M)^2R^2/\varepsilon^2\\{}+AR^2/\varepsilon+L_1^2R^2\end{gathered}\right)$.
\par Includes $\nabla f$, $h'$ and $\Pi_Q$ calls, plus $F$ values for selection; $\Pi_Q$ is unnecessary when $Q=E$.
& \textbf{(1)}: gradients $\varepsilon^{-1/2}$ instead of $\varepsilon^{-2}$ for fixed parameters, $G+M>0$; no $F$ values. Total gains can occur for constrained $G\gg M$. For $Q=E$, take $G=M$: matched total order unless $\varepsilon A\gg M^2$ or $\varepsilon L_1\gg M$; initialization can mask gains.
& \textbf{(1)}: updates additionally need $L_0,L_1,M,\Delta,\varepsilon$; \textbf{(3)} needs only $R,T$. Initialization term $(1+L_1R)\ell$ exceeds $L_1^2R^2$ when $\ell\gg L_1R\ge1$. \\
\bottomrule
\end{tabularx}
\begin{tablenotes}[flushleft]\footnotesize
\item[] \textbf{Abbr:} $A=L_0+L_1G$; $\ell=1+\log\!\left(1+(1+L_1R)\sqrt{\Delta/(L_0R^2)}\right)$; $\Pi_C$ is exact Euclidean projection onto $C$; $h'$ denotes a returned subgradient; $T$ is the iteration horizon in \textbf{(3)}. Total cost counts all primitive oracle calls.
\end{tablenotes}
\end{threeparttable}
\end{table}
\clearpage
